\section[Cote 67 : modules à involution et triples (P, Q, m)]{Cote 67 : modules à involution et triples $(P, Q, m)$\protect\verifie{Folder67}}\label{sec:67}

La cote 67, \emph{Chirurgie des surfaces conformes : notes manuscrites (s.d., 1983-1984), lettres (1984)}, cent pages que l'inventaire range dans le groupe \q{Autour de l'enseignement}, s'achève sur des notes consacrées aux courbes elliptiques réelles (pages~90 à~100). Une courbe elliptique réelle y revient à un réseau $\Pi$ de rang~$2$ muni d'une involution qui renverse l'orientation (page~92), et c'est pour classer ces involutions que la page~90, datée \q{Déc 83}, fait un détour d'algèbre linéaire sur un anneau quelconque, \q{par exemple} commutatif, le mot étant lu avec doute (\lot{67}{5}{90}). La page considère les $k$-modules $M$ munis d'un automorphisme $\sigma$ tel que $\sigma^2 = \mathrm{id}$ et sur lesquels la multiplication par~$2$ est injective, leur associe les triples $(P, Q, m)$, établit $M_+ \cap M_- = 0$, $2M \subset M_+ \oplus M_-$, $M' \cap P = 2P$ et $M' \cap Q = 2Q$ pour $M' = 2M$, reconstitue $(M, \sigma)$ à partir du triple par l'isomorphisme $2\,\mathrm{id}_M : M \to M'$, et énonce la proposition. La page~91 en est le brouillon ; la page~99 en tire le cas du rang~$2$ sur un anneau principal (\lot{67}{5}{99}).

\noindent\emph{Conventions.} Soit $k$ un anneau, non nécessairement commutatif. Un $k$-module $M$ est \emph{$2$-régulier} si $2x = 0$ entraîne $x = 0$ dans $M$ ; on note $2M$ l'image de la multiplication par~$2$, qui est un sous-module puisque $2$ est central. La catégorie $\mathcal I$ a pour objets les couples $(M, \sigma)$ d'un $k$-module $2$-régulier et d'un endomorphisme $k$-linéaire $\sigma$ tel que $\sigma^2 = \mathrm{id}$, et pour morphismes $(M, \sigma) \to (M', \sigma')$ les applications $k$-linéaires $f$ telles que $f\sigma = \sigma' f$. On pose $M_+ = \operatorname{Ker}(1 - \sigma)$ et $M_- = \operatorname{Ker}(1 + \sigma)$. La catégorie $\mathcal T$ a pour objets les triples $(P, Q, m)$ de deux $k$-modules $2$-réguliers et d'un sous-module $m$ de $P/2P \times Q/2Q$ tel que $(p, 0) \in m$ entraîne $p = 0$ et $(0, q) \in m$ entraîne $q = 0$ ; ses morphismes $(P, Q, m) \to (P', Q', m')$ sont les couples $(f, g)$ d'applications linéaires $P \to P'$, $Q \to Q'$ telles que $(\bar f \times \bar g)(m) \subset m'$, où $\bar f$, $\bar g$ sont les applications induites, qui existent puisque $f(2P) \subset 2P'$. Un sous-$k$-module de $P/2P \times Q/2Q$ est la même chose qu'un sous-$k/2k$-module, $2$ annulant ce produit. Pour un triple $T = (P, Q, m)$, on note $N_T \subset P \times Q$ l'image réciproque de $m$ : $(p, q) \in N_T$ si et seulement si $(\bar p, \bar q) \in m$.

\begin{enonce}[Page 90, \emph{proposition}\piste{67-involution-modules-triples}]
Soit $(M, \sigma)$ un objet de $\mathcal I$, et $P = M_+$, $Q = M_-$.
\begin{enumerate}
\item $P \cap Q = 0$ ; $x + \sigma x \in P$ et $x - \sigma x \in Q$ pour tout $x \in M$ ; $P \cap 2M = 2P$ et $Q \cap 2M = 2Q$.
\item Soit $m$ l'image dans $P/2P \times Q/2Q$ de $\{(p, q) \in P \times Q \mid p + q \in 2M\}$. Pour $p \in P$, $q \in Q$, on a $(\bar p, \bar q) \in m$ si et seulement si $p + q \in 2M$, et $(P, Q, m)$ est un objet de $\mathcal T$.
\item Le foncteur $B : \mathcal T \to \mathcal I$, $(P, Q, m) \mapsto (N_T, \mathrm{id}_P \times (-\mathrm{id}_Q))$, est une équivalence de catégories. L'application $x \mapsto (x + \sigma x, x - \sigma x)$ est un isomorphisme de $(M, \sigma)$ sur $B(P, Q, m)$ ; un quasi-inverse de $B$ envoie donc $(M, \sigma)$ sur un triple isomorphe à $(M_+, M_-, m)$, et le foncteur $(M, \sigma) \mapsto (P, Q, m)$ de la page est une équivalence.
\end{enumerate}
Aucune hypothèse de finitude ni de projectivité n'est faite sur $M$, $P$ ou $Q$.
\end{enonce}

\begin{proof}[Démonstration de (1)]
Si $x \in P \cap Q$, on a $\sigma x = x$ et $\sigma x = -x$, donc $2x = x + x = x - x = 0$, et $x = 0$ par $2$-régularité. Comme $\sigma^2 = \mathrm{id}$, $\sigma(x + \sigma x) = \sigma x + x$ et $\sigma(x - \sigma x) = -(x - \sigma x)$. Soit $p \in P$ avec $p = 2y$, $y \in M$ : alors $2\sigma y = \sigma(2y) = \sigma p = p = 2y$, donc $\sigma y = y$ par $2$-régularité, et $y \in P$, d'où $p \in 2P$. De même, si $q \in Q$ et $q = 2y$, alors $2\sigma y = -q = 2(-y)$, donc $\sigma y = -y$ et $y \in Q$. Les inclusions inverses $2P \subset P \cap 2M$, $2Q \subset Q \cap 2M$ sont claires.
\end{proof}

\begin{proof}[Démonstration de (2)]
Si $(\bar p, \bar q) \in m$, il existe $(p', q')$ avec $p' + q' \in 2M$, $p' - p \in 2P$ et $q' - q \in 2Q$ ; comme $2P$ et $2Q$ sont contenus dans $2M$, $p + q = (p' + q') - (p' - p) - (q' - q)$ est dans $2M$. La réciproque est la définition de $m$. Les modules $P$ et $Q$ sont $2$-réguliers comme sous-modules de $M$. Si $(\bar p, 0) \in m$, alors $(\bar p, \bar 0) \in m$, donc $p = p + 0 \in 2M$, donc $p \in 2P$ par~(1), c'est-à-dire $\bar p = 0$ ; de même pour $Q$.
\end{proof}

\begin{proof}[Démonstration de (3)]
\emph{Le foncteur.} Pour $T = (P, Q, m)$, l'application $\tau = \mathrm{id}_P \times (-\mathrm{id}_Q)$ conserve $N_T$, car $-\bar q = \bar q$ dans $Q/2Q$ (puisque $-q - q = 2(-q)$) ; on a $\tau^2 = \mathrm{id}$, et $N_T$ est $2$-régulier comme sous-module de $P \times Q$. Un morphisme $(f, g)$ envoie $N_T$ dans $N_{T'}$ par la condition $(\bar f \times \bar g)(m) \subset m'$, et $f \times g$ commute à $\tau$ ; c'est le foncteur $B$. Pour $p \in P$ et $q \in Q$, $(2p, 0)$ et $(0, 2q)$ sont dans $N_T$, leurs classes étant nulles.

\emph{Fidélité.} Si $B(f, g) = B(f', g')$, on évalue en $(2p, 0)$ : $2f(p) = 2f'(p)$, donc $f(p) = f'(p)$ par $2$-régularité de $P'$ ; de même $g = g'$ en évaluant en $(0, 2q)$.

\emph{Plénitude.} Soit $h : B(T) \to B(T')$ un morphisme de $\mathcal I$, $T' = (P', Q', m')$. Pour $p \in P$, posons $h(2p, 0) = (a, b)$. Comme $(2p, 0)$ est fixe par $\tau$, $(a, -b) = \tau h(2p, 0) = h(2p, 0) = (a, b)$, donc $2b = b + b = b - b = 0$ et $b = 0$. Comme $(a, 0) \in N_{T'}$, $(\bar a, 0) \in m'$, donc $\bar a = 0$ et $a \in 2P'$. La multiplication par~$2$ étant un isomorphisme linéaire de $P'$ sur $2P'$, on pose $f(p) = a/2$ ; $f$ est linéaire, comme composé de $p \mapsto h(2p, 0)$, de la première projection et de cet inverse. De même, $(0, 2q)$ étant changé de signe par $\tau$, $h(0, 2q) = (0, c)$ avec $c \in 2Q'$, et l'on pose $g(q) = c/2$. Pour $(p, q) \in N_T$, on a $2(p, q) = (2p, 0) + (0, 2q)$ dans $N_T$, donc
\[
2\,h(p, q) = h(2p, 0) + h(0, 2q) = (2f(p), 0) + (0, 2g(q)) = 2\,(f(p), g(q)),
\]
et $h(p, q) = (f(p), g(q))$ par $2$-régularité de $P' \times Q'$. En particulier, si $(\bar p, \bar q) \in m$, alors $(p, q) \in N_T$ et $(f(p), g(q)) = h(p, q) \in N_{T'}$, c'est-à-dire $(\overline{f(p)}, \overline{g(q)}) \in m'$ : $(f, g)$ est un morphisme de $\mathcal T$, et $B(f, g) = h$.

\emph{Essentielle surjectivité.} Soit $(M, \sigma)$ un objet de $\mathcal I$ et $T = (M_+, M_-, m)$ le triple de~(2). L'application $\varphi(x) = (x + \sigma x, x - \sigma x)$ est linéaire de $M$ dans $M_+ \times M_-$ par~(1), à valeurs dans $N_T$ par~(2), puisque $(x + \sigma x) + (x - \sigma x) = 2x$. Elle commute aux involutions : $\varphi(\sigma x) = (\sigma x + x, \sigma x - x) = \tau \varphi(x)$. Elle est injective : si $\varphi(x) = \varphi(y)$, alors $2x = (x + \sigma x) + (x - \sigma x) = 2y$, et $x = y$. Elle est surjective : si $(p, q) \in N_T$, alors $p + q \in 2M$ par~(2), soit $p + q = 2x$ ; on a $2\sigma x = \sigma p + \sigma q = p - q$, donc $2(x + \sigma x) = 2p$ et $2(x - \sigma x) = 2q$, d'où $\varphi(x) = (p, q)$. Un morphisme de $\mathcal I$ dont l'application sous-jacente est bijective est un isomorphisme, son inverse commutant encore aux involutions ; donc $(M, \sigma) \simeq B(T)$.

Un foncteur fidèle, plein et essentiellement surjectif est une équivalence. Si $C$ en est un quasi-inverse, $B(C(M, \sigma)) \simeq (M, \sigma) \simeq B(M_+, M_-, m)$, et $B$ étant pleinement fidèle, $C(M, \sigma) \simeq (M_+, M_-, m)$.
\end{proof}

\begin{remarque}
L'énoncé tient tel que la piste le formule, pour un anneau quelconque, commutatif ou non, et sans hypothèse de finitude ni de projectivité. La seule hypothèse qui serve est la $2$-régularité, à trois endroits : $P \cap Q = 0$ ; $P \cap 2M = 2P$ et $Q \cap 2M = 2Q$, qui font que $m$ rencontre trivialement chaque facteur ; et la plénitude, où un morphisme $h$ est déterminé par sa restriction à $2P \times 2Q$ parce que $2N_T \subset 2P \times 2Q$ et que le but est $2$-régulier. La démonstration suit la page : la reconstruction de $(M, \sigma)$ y est exactement l'isomorphisme $\varphi$, qui est la multiplication par~$2$ de $M$ sur $M' = 2M$ lue dans $P \oplus Q$.
\end{remarque}

\begin{remarque}
La vérification énonce l'équivalence dans le sens $B : \mathcal T \to \mathcal I$, de la reconstruction, et en déduit le foncteur de la page, $(M, \sigma) \mapsto (P, Q, m)$, comme quasi-inverse, à isomorphisme près sur chaque objet. Les morphismes de triples ne sont pas définis sur la page ; on a pris les couples $(f, g)$ tels que $(\bar f \times \bar g)(m) \subset m'$, qui sont ceux pour lesquels $f \times g$ envoie $M'$ dans $M'$. La condition sur $m$, que la page écrit $m \cap P/2P = 0$, $m \cap Q/2Q = 0$, est prise telle quelle ; sa traduction en graphe d'un isomorphisme entre un sous-module de $P/2P$ et un sous-module de $Q/2Q$ n'est pas formalisée.
\end{remarque}

\noindent Ne sont pas démontrés ici : la description de $m$ comme graphe d'un isomorphisme $p' \simeq q'$ entre sous-modules de $P/2P$ et de $Q/2Q$ ; le premier corollaire de la page~90, pour $M$ projectif de rang~$2$ sur un anneau connexe, en partie illisible ; le corollaire de la page~99 (sur un anneau principal $2$-régulier avec $k/2k$ un corps, deux classes de conjugaison d'involutions distinctes de $\pm 1$ dans $\mathrm{GL}_2(k)$, représentées par $\left(\begin{smallmatrix} 1 & 0 \\ 0 & -1 \end{smallmatrix}\right)$ et $\left(\begin{smallmatrix} 1 & 1 \\ 0 & -1 \end{smallmatrix}\right)$), et son usage pour les deux types de réseaux des courbes elliptiques réelles (page~92) ; tout le reste du dossier, la chirurgie des surfaces conformes et les métriques canoniques.
